\newtoggle{withbonus} \toggletrue{withbonus} % set true to show bonus points in grade table

% setting underlying course name (Grundlagenfach Mathematik, §-Unterricht, \ldots)
\newcommand{\mycourse}{Ergänzungsfach Informatik}
% name of the class
\newcommand{\myclass}{4. Klassen}
% set date
\newcommand{\mydate}{24. Januar 2024}

\title{\textbf{Endliche Automaten}}
\author{Autoren: Thomas Graf\\
\mycourse, \myclass}
\date{\mydate}
\maketitle
\thispagestyle{headandfoot}		

\begin{center}
	\fbox{\parbox{140mm}{\centering
			\begin{itemize}
				\item Dauer der Prüfung: $40$ Minuten.
				%\item \textbf{Zeige in jeder Aufgabe unbedingt Deine Schritte!}
				\item Erlaubte Hilfsmittel:
				\begin{itemize}
				    \item Schreibzeug
				    %\item Taschenrechner
				\end{itemize}
	\end{itemize}}}
\end{center}
\vspace{10mm}

\begin{center}
	Vorname: \rule{45mm}{0.8pt}\qquad Nachname: \rule{45mm}{0.8pt}
\end{center}
\vspace{20mm}

\begin{center}
	\iftoggle{withbonus} {
		\combinedgradetable[h][questions]
	} {
		\gradetable[h][questions]
	}
\end{center}
\vspace{15mm}
\begin{center}
	Note: 
	\vspace{10mm}
	
	\rule{8mm}{1.2pt}
\end{center}
%\thispagestyle{empty}
\clearpage

\begin{questions}
\question{\textbf{Endlichen Automaten zeichnen 1}}

Gegeben ist die Sprache
\begin{align*}
  L_1 := \setcm{x00110y}{x,y\in\lrc{0,1}^*}.  
\end{align*}
\begin{parts}
  \part[1] Nenne ein kürzestes Wort in $L_1$.
  \vspace{3cm}
  \part[3] Entwerfe einen endlichen Automaten (in Diagrammform), welcher $L_1$ akzeptiert.
  \vspace{6cm}
\end{parts}
\droptotalpoints

\begin{questions}
\question{\textbf{Endlichen Automaten zeichnen 2}}

Gegeben ist die Sprache
\begin{align*}
  L_2 := \setcm{x1y}{x,y\in\lrc{0,1}^*, \abs{y}\geq 1}.  
\end{align*}
\begin{parts}
  \part[1] Beschreibe die Menge $L_2$ in präzisen Worten. 
  \vspace{3cm}
  \part[1] Nenne ein kürzestes Wort in $L_2$.
  \vspace{3cm}
  \part[3] Entwerfe einen endlichen Automaten (in Diagrammform), welcher $L_1$ akzeptiert.
  \vspace{6cm}
\end{parts}
\droptotalpoints


\question{\textbf{Nicht reguläre Sprache}}

Gegeben ist die Sprache
\begin{align*}
  L_{\text{equal}} := \setcm{w\#w}{w\in\lrc{a,b}^*}.  
\end{align*}
\begin{parts}
  \part[1] Nenne genau drei verschiedene Wörter, die in $L_{\text{equal}}$ liegen. 
  \vspace{3cm}
  \part[3] Zeige, dass $L_{\text{equal}}$ nicht regulär ist.
\end{parts}
\droptotalpoints


\end{questions}
